% Auto-generated from report/2.md — do not edit by hand.
\section{مبانی نظری یادگیری بازنمایی روی گراف}\label{sec:2-1}


\subsection{تعاریف پایه و نمادگذاری ریاضی}\label{sec:2-1-1}


یک شبکه تعاملی زیستی به‌صورت گراف بدون جهت $G = (\mathcal{V}, \mathcal{E})$ مدل‌سازی می‌شود:

\begin{itemize}
\item \textbf{مجموعه گره‌ها} $\mathcal{V}$: موجودیت‌های بیولوژیکی با کاردینالیتی $|\mathcal{V}| = N$.
\item \textbf{مجموعه یال‌ها} $\mathcal{E}$: برهم‌کنش‌های تجربی مشاهده‌شده با تعداد $|\mathcal{E}| = M$.
\item \textbf{ماتریس مجاورت} $A \in \{0, 1\}^{N \times N}$: ماتریسی متقارن که $A_{uv} = 1$ در صورت وجود یال $(u, v) \in \mathcal{E}$.
\item \textbf{ماتریس درجه} $D \in \mathbb{R}^{N \times N}$: ماتریس قطری با عناصر $D_{uu} = d_u = \sum_{v} A_{uv}$.
\item \textbf{ماتریس ویژگی‌های ورودی} $X \in \mathbb{R}^{N \times d_{in}}$: در این پژوهش، $X = I_N$ (ماتریس هویت تنک) انتخاب شده است (بخش ۱-۵-۲).
\end{itemize}

\subsection{پارادایم انتشار پیام (\lr{Message Passing})}\label{sec:2-1-2}


اکثریت قاطع شبکه‌های عصبی گراف مدرن از چارچوب شبکه‌های عصبی پیام‌رسان\footnote{شبکه‌های عصبی پیام‌رسان (\lr{Message Passing Neural Networks}).}\cite{gilmer2017mpnn} پیروی می‌کنند. در این رویکرد، بازنمایی هر گره $v$ در لایه $(l+1)$ از طریق چرخه «تجمیع» و «به‌روزرسانی» محاسبه می‌شود:


\begin{equation}
\label{eq:2.1}
m_v^{(l+1)} = \text{\lr{AGG}}^{(l+1)}\left( \left\{ \text{\lr{MSG}}^{(l+1)}\left(h_u^{(l)}, h_v^{(l)}\right) : u \in \mathcal{N}(v) \right\} \right)
\end{equation}



\begin{equation}
\label{eq:2.2}
h_v^{(l+1)} = \text{\lr{UPDATE}}^{(l+1)}\left( h_v^{(l)}, m_v^{(l+1)} \right)
\end{equation}


\subsection{شبکه‌های پیچشی گراف و نرمال‌سازی لاپلاسی}\label{sec:2-1-3}


کیپف و ولینگ\cite{kipf2017gcn} معماری \lr{GCN} را بر پایه نرمال‌سازی متقارن لاپلاسی فرموله کردند:


\begin{equation}
\label{eq:2.3}
\tilde{A} = A + I_N, \qquad \tilde{D}_{uu} = d_u + 1
\end{equation}



\begin{equation}
\label{eq:2.4}
F = \tilde{D}^{-1/2} \tilde{A} \tilde{D}^{-1/2}
\end{equation}



\begin{equation}
\label{eq:2.5}
H^{(l+1)} = \sigma\left( F H^{(l)} W^{(l)} \right)
\end{equation}


در پیاده‌سازی این پژوهش (تابع \pcode{laplacian\_normalize} در \pcode{src/data/normalization.py})، برای گره‌های منزوی به‌جای افزودن $\epsilon$ مصنوعی، مقدار صفر منظور می‌شود:


\begin{equation}
\label{eq:2.6}
(D^{-1/2})_{ii} = \begin{cases} d_i^{-0.5} & \text{\lr{if }} d_i > 0 \\ 0 & \text{\lr{otherwise}} \end{cases}
\end{equation}


این انتخاب از تقویت مصنوعی گره‌های منزوی جلوگیری می‌کند و در آزمون واحد \pcode{test\_laplacian\_isolated\_nodes\_are\_zero} اعتبارسنجی شده است.

\subsection{چالش بیش‌هموارسازی}\label{sec:2-1-4}


با افزایش عمق لایه‌ها، عملگر انتشار $F$ به‌عنوان فیلتر پایین‌گذر عمل کرده و بازنمایی گره‌ها به سمت همگرایی میل می‌کند\footnote{بیش‌هموارسازی (\lr{Over-smoothing}).}\cite{li2018deeper}:


\begin{equation}
\label{eq:2.7}
\lim_{l \to \infty} H^{(l)} \propto \mathbf{1} \cdot \pi^T
\end{equation}


در شبکه‌های زیستی تنک با قطر کوچک\footnote{خاصیت جهان‌کوچک (\lr{Small-World Property}).}\cite{watts1998smallworld}، این پدیده تشدید می‌شود. این چالش یکی از انگیزه‌های طراحی عملگرهای پرش توپولوژیک در این پژوهش است.

\section{پیش‌بینی پیوند و خطوط مبنای هوریستیک}\label{sec:2-2}


\subsection{صورت‌بندی مسئله}\label{sec:2-2-1}


اگر امبدینگ نهایی گره‌ها $E = H^{(L)} \in \mathbb{R}^{N \times d}$ باشد، احتمال یال میان $u$ و $v$ توسط دیکودر تخمین زده می‌شود:


\begin{equation}
\label{eq:2.8}
\hat{y}_{uv} = \text{\lr{Decoder}}(E_u, E_v)
\end{equation}


\subsection{شاخص‌های هوریستیک غیرپارامتری}\label{sec:2-2-2}


پیش از \lr{GNN}ها، پیش‌بینی پیوند با شاخص‌های ساختاری انجام می‌گرفت. جدول ۲-۱ چهار شاخص اصلی را خلاصه می‌کند:


\begin{table}[htbp]
\centering
\footnotesize
\caption{فرمولاسیون شاخص‌های هوریستیک کلاسیک}
\label{tab:2-1}
\begin{tabular}{lll}
\toprule
شاخص & فرمول & ویژگی \\
\midrule
همسایگان مشترک (\lr{CN}) & $|\mathcal{N}(u) \cap \mathcal{N}(v)|$ & شمارش ساده بدون جریمه درجه \\
ضریب جاکارد & $\frac{|\mathcal{N}(u) \cap \mathcal{N}(v)|}{|\mathcal{N}(u) \cup \mathcal{N}(v)|}$ & نرمال‌سازی با اجتماع \\
آدامیک-آدار (\lr{AA}) & $\sum_{k} \frac{1}{\log(\max(2, d_k))}$ & جریمه لگاریتمی هاب‌ها \\
تخصیص منابع (\lr{RA}) & $\sum_{k} \frac{1}{\max(1, d_k)}$ & جریمه خطی هاب‌ها \\
\bottomrule
\end{tabular}
\end{table}


شاخص تخصیص منابع\cite{zhou2009predicting} در این پژوهش هم به‌عنوان عملگر پرش پیوسته (بخش ۳-۳-۲) و هم به‌عنوان خط مبنای غیرپارامتری (\pcode{src/models/heuristics.py}) به‌کار رفته است. در گراف‌های همگن، امتیاز بر پایه همسایگان مشترک وزن‌دار محاسبه می‌شود. در گراف‌های دوقطبی، از آنجا که همسایگان مشترک ۱-گام تهی هستند، از ضرب ماتریسی ۳-گام استفاده می‌شود:


\begin{equation}
\label{eq:2.9}
\text{\lr{Score}}(u, v) = \left( A D^{-1} A^T D^{-1} A \right)_{uv}
\end{equation}


این فرمولاسیون در تابع \pcode{\_bipartite\_three\_walk\_scores} پیاده‌سازی شده و در آزمون \pcode{test\_bipartite\_heuristics\_are\_nonzero} صحت‌سنجی گردیده است.

\section{تشریح مدل مرجع \lr{SkipGNN}}\label{sec:2-3}


مدل \lr{SkipGNN}\cite{huang2020skipgnn} با فرض «شباهت پرشی» معرفی شد: موجودیت‌هایی که مستقیماً متصل نیستند اما همسایگان مشترک دارند، شباهت ساختاری دارند.

\subsection{ساخت گراف پرش باینری}\label{sec:2-3-1}



\begin{equation}
\label{eq:2.10}
A_s = \text{\lr{sign}}(A A^T)
\end{equation}


\subsection{معماری ادغام دوکاناله}\label{sec:2-3-2}


لایه اول انکودر:


\begin{equation}
\label{eq:2.11}
H^{(1)} = \text{\lr{ReLU}}\left( F X W_o + F_s X W_{so} \right)
\end{equation}



\begin{equation}
\label{eq:2.12}
S^{(1)} = \text{\lr{ReLU}}\left( F_s X W_s + F H^{(1)} W_{os} \right)
\end{equation}


لایه دوم و تجمیع نهایی:


\begin{equation}
\label{eq:2.13}
E = F \hat{H}^{(1)} W_o^{(1)} + F_s \hat{S}^{(1)} W_{s2o}^{(1)}
\end{equation}


\subsection{دیکودر خطی}\label{sec:2-3-3}



\begin{equation}
\label{eq:2.14}
\text{\lr{feat}}_{uv} = [E_u \parallel E_v] \in \mathbb{R}^{2d_2}
\end{equation}



\begin{equation}
\label{eq:2.15}
\text{\lr{Logit}}_{uv} = W_{d2} \cdot \text{\lr{ReLU}}(W_{d1} \text{\lr{feat}}_{uv} + b_{d1}) + b_{d2}
\end{equation}


این مدل در کد این پژوهش به‌صورت \pcode{SkipGNNBaseline} (\pcode{src/models/skipgnn.py}) بازپیاده‌سازی شده است.

\section{نقد ساختاری و ریاضی \lr{SkipGNN}}\label{sec:2-4}


تحلیل ریاضی و تجربی در این پژوهش، سه محدودیت بنیادین در معماری \lr{SkipGNN} شناسایی کرده است.

\subsection{بن‌بست هندسی در گراف‌های دوقطبی (اثبات ریاضی)}\label{sec:2-4-1}


\textbf{قضیه:} در گراف دوقطبی با افراز $\mathcal{V} = \mathcal{V}_1 \cup \mathcal{V}_2$، ماتریس $A^2$ ساختاری بلوکی-قطری دارد و بلوک‌های غیرقطری آن صفر هستند.

\textbf{اثبات:} ماتریس مجاورت گراف دوقطبی به فرم بلوکی زیر است:


\begin{equation}
\label{eq:2.16}
A = \begin{bmatrix} \mathbf{0} & B \\ B^T & \mathbf{0} \end{bmatrix}
\end{equation}


محاسبه $A^2$:


\begin{equation}
\label{eq:2.17}
A^2 = \begin{bmatrix} \mathbf{0} & B \\ B^T & \mathbf{0} \end{bmatrix} \begin{bmatrix} \mathbf{0} & B \\ B^T & \mathbf{0} \end{bmatrix} = \begin{bmatrix} BB^T & \mathbf{0} \\ \mathbf{0} & B^TB \end{bmatrix}
\end{equation}


بلوک‌های غیرقطری صفر مطلق هستند. بنابراین، گراف پرش $A_s = \text{\lr{sign}}(A^2)$ هیچ یالی میان $\mathcal{V}_1$ و $\mathcal{V}_2$ ایجاد نمی‌کند. $\blacksquare$

\textbf{نتیجه:} در مسئله پیش‌بینی \lr{DTI}، کانال پرش \lr{SkipGNN} پیام‌ها را فقط میان دارو-دارو و پروتئین-پروتئین منتقل می‌کند و اطلاعاتی درباره تارگت‌های کاندید از نوع مخالف ارائه نمی‌دهد.

\textbf{مسیر ۳-گام به‌عنوان اولین پل:} محاسبه $A^3$ نشان می‌دهد که بلوک‌های غیرقطری فعال می‌شوند:


\begin{equation}
\label{eq:2.18}
A^3 = \begin{bmatrix} \mathbf{0} & BB^TB \\ B^TBB^T & \mathbf{0} \end{bmatrix}
\end{equation}


بنابراین، مسیر $d \to p' \to d' \to p$ (دارو \faarrow{} پروتئین واسط \faarrow{} داروی واسط \faarrow{} پروتئین هدف) اولین مسیر توپولوژیک غیرمستقیم برای دسترسی به تارگت کاندید است. این بینش، انگیزه طراحی عملگر ۳-گام در مدل‌های \pcode{ThreeHopSkipGNN} و هوریستیک دوقطبی است.

\subsection{حذف اطلاعات چگالی بر اثر باینری‌سازی}\label{sec:2-4-2}


عملگر $\text{\lr{sign}}$ در معادله (۲-۱۰)، تمایز میان جفت‌هایی با تعداد متفاوت همسایه مشترک را حذف می‌کند:


\begin{equation}
\label{eq:2.19}
(A_s)_{uv} = \text{\lr{sign}}\left(|\mathcal{N}(u) \cap \mathcal{N}(v)|\right) \in \{0, 1\}
\end{equation}


جفتی با ۵۰ همسایه مشترک و جفتی با ۱ همسایه مشترک، وزن یکسانی دریافت می‌کنند. در گراف متراکم \lr{DDI} (میانگین درجه ۶۴)، این باینری‌سازی گراف پرش را به گرافی نزدیک به کامل تبدیل می‌کند.

\textbf{راه‌حل پیشنهادی:} عملگر تخصیص منابع (بخش ۳-۳-۲):


\begin{equation}
\label{eq:2.20}
W_2 = A D^{-1} A^T, \qquad (W_2)_{uv} = \sum_{k \in \mathcal{N}(u) \cap \mathcal{N}(v)} \frac{1}{d_k}
\end{equation}


\subsection{تنگنای ظرفیت دیکودر خطی}\label{sec:2-4-3}


دیکودر معادله (۲-۱۵) فاقد مؤلفه‌های تعاملی ضربی و فاصله‌ای است. در گراف‌های بدون‌جهت، تقارن پیش‌بینی ($\hat{y}_{uv} = \hat{y}_{vu}$) تضمین نمی‌شود.

\textbf{راه‌حل پیشنهادی:} دیکودر تانسوری ۴‌گانه (بخش ۴-۲-۴):


\begin{equation}
\label{eq:2.21}
Z_{uv} = [E_u \parallel E_v \parallel E_u \odot E_v \parallel |E_u - E_v|]
\end{equation}


\section{نقد روش‌شناختی پروتکل ارزیابی}\label{sec:2-5}


\subsection{سوگیری درجه در نمونه‌برداری منفی یکنواخت}\label{sec:2-5-1}


در پروتکل ارزیابی مقاله مرجع، منفی‌ها به‌صورت تصادفی یکنواخت انتخاب می‌شدند. با توجه به توزیع دم‌سنگین درجه در گراف‌های زیستی\cite{barabasi1999}، بخش عمده منفی‌های تصادفی از گره‌های کم‌درجه تشکیل می‌شوند. شواهد تجربی این پژوهش (جدول~\ref{tab:5-3}) نشان می‌دهد که عملکرد \lr{SkipGNN} در آزمون یکنواخت ($0.929$ در \lr{DTI}) به‌طور قابل‌توجهی بالاتر از آزمون با منفی‌های هم‌درجه ($0.695$) است.

این یافته با نقد لی و همکاران\cite{li2023evaluating} درباره محدودیت‌های بنچمارک‌های مبتنی بر منفی‌های تصادفی هم‌راستا است.

\subsection{ناسازگاری‌های پیاده‌سازی در کد مرجع}\label{sec:2-5-2}


بازبینی کد رسمی \lr{SkipGNN}\footnote{مخزن کد مرجع: \pcode{github.com/kexinhuang12345/SkipGNN}}\cite{huang2020skipgnncode}، چند ناسازگاری نرم‌افزاری شناسایی کرد که در این پژوهش اصلاح شدند:

\textbf{ناسازگاری اول  -  آستانه‌گذاری روی لاجیت خام:} در تابع ارزیابی کد مرجع، شرط $\text{\lr{Logit}} \ge 0.5$ اعمال می‌شد، معادل احتمال $\sigma(0.5) \approx 0.622$. این پروژه آستانه $\tau^*$ را صرفاً روی اعتبارسنجی و با پویش ۹۱ نقطه‌ای محاسبه می‌کند (\pcode{src/eval/metrics.py}).

\textbf{ناسازگاری دوم  -  فرمت تنک/متراکم:} در لودر \lr{DDI} کد مرجع، ماتریس مجاورت به فرمت متراکم تبدیل می‌شد، در حالی که لایه \lr{GCN} انتظار تنسور تنک داشت. این پروژه تمام ماتریس‌ها را در فرمت \pcode{torch.sparse\_coo\_tensor} کوالس‌شده نگه می‌دارد.

\textbf{ناسازگاری سوم  -  نرمال‌سازی ویژگی‌ها:} نرمال‌سازی سطری $L_1$ روی امبدینگ‌های \lr{Node2Vec} (که شامل مقادیر منفی هستند) می‌توانست اعوجاج ایجاد کند. این پروژه از استانداردسازی ستونی \lr{Z-Score} استفاده می‌کند.

\section{جمع‌بندی: شکاف پژوهشی و ضرورت مدل‌های پیشنهادی}\label{sec:2-6}


جدول ۲-۲ خلاصه‌ای از محدودیت‌های شناسایی‌شده و راه‌حل‌های ارائه‌شده در این پژوهش را نشان می‌دهد:


\begin{table}[htbp]
\centering
\footnotesize
\caption{مقایسه محدودیت‌های \lr{SkipGNN} و راه‌حل‌های پیشنهادی}
\label{tab:2-2}
\begin{tabularx}{\textwidth}{>{\hspace{0pt}}p{0.22\textwidth}>{\hspace{0pt}}X>{\hspace{0pt}}X}
\toprule
مؤلفه & وضعیت در \lr{SkipGNN} & راه‌حل در این پژوهش \\
\midrule
عملگر پرش & باینری $\text{\lr{sign}}(AA^T)$ & پیوسته تخصیص منابع $AD^{-1}A^T$ \\
گراف دوقطبی & بن‌بست گام زوج $A^2$ & مسیرهای ۳-گام ($W_2D^{-1}A$) \\
ادغام لایه‌ای & جمع ثابت & گیتینگ تطبیقی سیگموئیدی \\
دیکودر & الحاق خطی $[E_u \parallel E_v]$ & تانسوری ۴‌گانه + \lr{BatchNorm} \\
پایداری فضا & فاقد قید & زیان متقارن \lr{InfoNCE} \\
پروتکل ارزیابی & یکنواخت & دوگانه (یکنواخت + سخت هم‌درجه) \\
\bottomrule
\end{tabularx}
\end{table}


فصول بعدی به تشریح داده‌ها (فصل ۳)، فرمولاسیون ریاضی چهار مدل پیشنهادی (فصل ۴)، نتایج تجربی (فصل ۵) و بحث (فصل ۶) اختصاص دارند.
